\pg{442} a) E' a destruição da \emph{sciencia}. Na verdade, a sciencia é o conhecimento das \emph{causas} dos seres. Ora, se \textsc{Deus} não existisse, os problemas mais fundamentaes seriam inexplicaveis. Porquanto, --- na ordem \emph{physica}, seria inexplicavel a origem do mundo, da vida, da alma humana, a uniformidade e constancia da ordem e das leis physicas; --- na ordem \emph{logica}, faltaria a base da \emph{possibilidade} das coisas, da \emph{eternidade} e \emph{necessidade} das essencias, da \emph{immutabilidade} da verdade; --- na ordem \emph{moral}, seria inexplicavel a existencia da \emph{lei moral}, a virtude não teria o premio, nem o vicio o castigo, e o bem não differiria do mal.

b) E' a destruição da \emph{sociedade}. Porquanto, uma sociedade não póde subsistir, sem que haja fidelidade e respeito nos subditos, desinteresse e justiça no principe, amor e concordia em todos. Mas uma sociedade sem \textsc{Deus} é o reino da anarchia, tendo por direito a força material, e por lei o despotismo \footnote{\emph{Voltaire} diz (\emph{Dict. phil.}, a. \emph{athée}): « Eu não quereria ser subdito de um principe atheu, que, se encontrasse o seu interesse em me esmigalhar, com certeza esmigalhar-me-hia. Eu não quereria, se fosse soberano, ter cortezãos atheus, pois, se tivessem interesse em me envenenar, deveria tomar, todos os dias, contra-venenos. E', pois, absolutamente necessario, para os principes e para os povos, que a idea de um Ente supremo, creador, governador, remunerador, seja gravada profundamente em todos os espiritos ».}.

c) E' a destruição da \emph{natureza humana}. Com effeito, se \textsc{Deus} não existisse, a nossa alma não seria immortal, e assim o homem seria a mais desgraçada das creaturas; porque, depois de ter soffrido tantas dores, depois de ter suspirado tanto pela felicidade sem nunca a encontrar, o \emph{nada} seria a recompensa dos seus soffrimentos e o descanço das suas tendencias \footnote{Estas consequencias do \emph{atheismo} são tambem provas, ainda que \emph{indirectas}, da nossa these. Comtudo, \emph{Stuart Mill} diz que não temos prova alguma, nem \emph{seriamente provavel}, para demonstrar a existencia de \textsc{Deus}!}.
